\documentclass{article}
\title{MAT1105 Oblig 3}
\author{borgebj}
\date{}

\usepackage[norsk]{babel}
\usepackage{subcaption}
\usepackage{amsmath}
\usepackage{amssymb}
\usepackage{graphicx}
\usepackage[a4paper, top=1in, bottom=1in, left=1in, right=1in]{geometry}
\usepackage{subcaption}
\usepackage{float}
\setlength{\parskip}{2em}
\setlength{\parindent}{0pt}
\linespread{1}


\begin{document}
\maketitle

\section*{Seksjon 1.9}

\subsection*{Oppgave 2}

\(T : \mathbb{R}^2 \rightarrow \mathbb{R}^4\) som tilfredsstiller\\ \\
\(T(e_1) = \begin{pmatrix}
    -1 \\ 2 \\ -3 \\ 4
\end{pmatrix}\) \ og \ \(T(e_2) = \begin{pmatrix}
    0 \\ -2 \\ 4 \\ 7
\end{pmatrix}\)
\\ \\ \\ 
Funksjonen tar inn en vektor i planet og gir ut en vektor i det 4-dimensjonelle rommet. \\ 
Detvil si at $e_1$ og $e_2$ er første og andre søyle i matrisen A.\\ \\
Da er matrisen til T, \quad \(\underline{\underline{A = 
\begin{pmatrix}
    -1 & 0 \\ 2 & -2 \\ -3 & 4 \\ 4 & 7
\end{pmatrix}}}\)

\newpage

\subsection*{Oppgave 3}
\say{\(a, b \in \mathbb{R}^2, \quad T : \mathbb{R}^2 \rightarrow \mathbb{R}^2\) tilfredsstiller \(\quad 
T(a) = \begin{pmatrix}
    -2 \\ 1
\end{pmatrix}\) \ og \ 
\(T(b) = \begin{pmatrix}
    0 \\ 3
\end{pmatrix}\)}

Egenskapen av en lineæravbildning er at:
\begin{enumerate}
    \item \(T(a+b) = T(a) + T(b)\)
    \item \(T(\textbf{c}x) = \textbf{c}T(x)\), for $c \in \mathbb{R}$
\end{enumerate}

\(T(3a-2b) \text{ kan skrives som: } \\\\ T(3a + (-2b)) \quad=\quad T(3a) + T(-2b) \quad=\quad 3T(a) + (-2)T(b) \quad=\quad 3T(x) - 2T(b)\)

\(3T(a) = 3\cdot\begin{pmatrix}
    -2 \\ 1
\end{pmatrix} = \begin{pmatrix}
    -6 \\ 3
\end{pmatrix}
\quad \quad \quad
-2T(b) = -2\cdot\begin{pmatrix}
    0 \\ 3
\end{pmatrix} = \begin{pmatrix}
    0 \\ -6
\end{pmatrix}
\quad \Leftrightarrow\\ \\ \\ \Leftrightarrow \quad
3T(x)-2T(b) \quad=\quad \begin{pmatrix}
    -6 \\ 3
\end{pmatrix} + \begin{pmatrix}
    0 \\ -6
\end{pmatrix} \quad=\quad \underline{\underline{\begin{pmatrix}
    -6 \\ -3
\end{pmatrix} \ = \ T(3a-2b)}}\)

\end{document}
